\section{Related work}
\label{sec:related}

\paragraph{Sparse and hierarchical attention.} Fixed sparsity patterns~\cite{child2019sparse,beltagy2020longformer,
zaheer2020bigbird} and content-based routing~\cite{kitaev2020reformer,roy2021routing} reduce attention's quadratic
cost; Routing Transformer already reached $O(n^{1.5})$ with content-based clustering~\cite{roy2021routing}.
Block-level selection is now common: Quest chooses key pages per query at inference time~\cite{tang2024quest}, MoBA
routes each query to a few trained blocks~\cite{lu2025moba}, Landmark Attention and InfLLM retrieve blocks through
representative tokens~\cite{mohtashami2023landmark,xiao2024infllm}, DeepSeek Sparse Attention selects tokens with a
lightweight learned indexer~\cite{deepseek2025v32}, and NSA~\cite{yuan2025nsa} combines compressed, selected and windowed
branches with trainable sparsity; our selector inherits NSA's three-branch layout. Hierarchical selection is well established: HiP prunes a tree of key blocks without
training~\cite{lee2025hip}, HISA filters pooled blocks before refining to tokens~\cite{xu2026hisa}, gist-token methods
learn chunk summaries and recursive summary hierarchies~\cite{mao2026gist}, and HiLS trains chunk-level retrieval scores
through the language-model loss~\cite{hu2026hils}. Several of these papers vary their selection budget and report
retrieval-sensitive benchmarks. Our
selector is a simpler member of this family, and we claim no algorithmic novelty. HiLS is closest in spirit; it
measures retrieval in small models under matched sparse budgets on RULER. We differ in scale and in protocol: one
synthetic task, but a dense model on byte-identical splits for every comparison, reserved confirmation seeds, and
every failed run reported. The Sparse Frontier studies sparsity--accuracy trade-offs of training-free methods in
large pretrained models across tasks and lengths~\cite{nawrot2025sparsefrontier}; our setting is the opposite end:
models trained from scratch with the selector, at toy scale. Two recent methods train selection end to end: SAS injects
the selector's continuous scores into attention logits so the language-modelling loss trains the ranking, and
compares gate forms~\cite{li2026sas}; FreeSolo pairs a sliding window with a learned long-range cell selector and
reports zero-shot length extrapolation~\cite{shi2026freesolo}. Our failed router gate (Section~\ref{sec:fixes}) is one
implementation on one task and does not test those designs.

\paragraph{Recall in efficient sequence models.} Zoology introduced MQAR and showed that most of the quality gap
between attention and gated-convolution models comes from associative recall~\cite{arora2024zoology}. Later work
traced a trade-off between recall and recurrent state size in linear attention~\cite{arora2024based} and proved and measured that fixed-state
models copy worse than transformers~\cite{jelassi2024repeat}; MQAR results for recurrent models are also sensitive to the
learning rate, which can make capacity conclusions from untuned runs misleading~\cite{okpekpe2025revisiting}. Long-context evaluations show that models use
information unevenly across their inputs~\cite{liu2024lost}. These studies compare architectures; we hold the
architecture fixed and vary how much of the context a sparse selector may consider.

\paragraph{Efficiency claims and their evaluation.} Kernel-level work such as FlashAttention reports speed with its
hardware, lengths and baselines stated~\cite{dao2022flashattention}. Some recent commercial reports pair large
long-context speed-ups with benchmark scores but without a dense model trained under matched
settings~\cite{subq2026,appen2026}. We do not evaluate those systems, whose mechanisms are not public. We use identical data to remove one source of confounding;
the remaining mixer and optimisation differences (Section~\ref{sec:dense}) still prevent attributing the whole
dense--selector gap to sparsity alone, and within-model controls are an equally valid design.
